\section{Evaluation}
\label{sec:eval}

\subsection{Random Benchmarks}~\label{sec:eval:random}

\subsection{Synthetic Benchmarks}~\label{sec:eval:synthetic}


\subsection{Hardware Benchmarks}~\label{sec:eval:hardware}

In this section, we evaluate on a set of benchmarks checking the equivalence of
two circuits. The canonical way to check circuit equivalence is using a miter, a
construction where given two circuits $C_1$ and $C_2$ (with the same number of
inputs and outputs), we construct a new circuit $M$ that (1) for each output of
$C_1$ and $C_2$, takes their XOR and (2) ORs all the XOR results together. We
then check whether the output of $M$ is always 0, which indicates that $C_1$ and
$C_2$ are equivalent.

Typically, the miter is converted into conjunctive normal form (CNF) and given to 
a SAT solver. However, these formulas can be very large and difficult to solve.
Biere et al.~\cite{clausal-cc} showed that often different circuits have redundant subcircuits and 
thus these formulas can be simplified significantly using congruence closure.

Biere et al. evaluated their clausal congruence closure algorithm on two benchmark sets:
five circuit pairs that curated to be difficult for SAT solvers~\cite{iwls} and a set of 
benchmarks comparing an original circuit against an ABC-optimized version~\cite{hwmcc}.
They implemented their congruence closure algorithm as a preprocessing pass in the SAT solver
Kissat~\cite{kissat} and showed that it significantly reduced the time taken to
solve these benchmarks. For our evaluation, we take the same benchmarks and run
them through our parallel congruence closure algorithms, comparing the time
taken to solve the benchmarks with and without congruence closure.


% \subsection{Equality Saturation Benchmarks}

% To get equality saturation benchmarks, we ran egglog~\cite{egglog} on the egg-cc
% project~\cite{egg-cc}—an experimental optimizing compiler for the LLVM-like IR 
% Bril. We ran egglog on a number of Bril programs for \as{todo} seconds and logged
% all of the unions egglog generated during this time. We then extracted these unions
% into and MST-LIB file and evaluated our parallel congruence closure algorithms on these
% benchmarks.

% \subsection{SMT Benchmarks}

% To get SMT benchmarks, we took benchmarks from the QF-UF logic
% in the annual 2025 SMT competition~\cite{smtcomp2025}. In this category, we picked
% the benchmarks from \asm{fill in} as the other benchmarks were either very simple
% or required more complicated boolean reasoning.

% These benchmarks include boolean connectives, so to get pure congruence closure
% constraints, we ran each of the bencmarks through the Sundance-SMT solver
% ~\cite{sundance}, every time Sundance-SMT found a conflict in its congruence 
% closure solver, we extracted the set of equalities into an SMT-LIB file.





